
\section{The global trace formula in stable(r) form}

Our primary references for this section are \cite{LL}, \cite{L3}.
\subsection{Definitions and a general discussion of the results}

We define $G_1=GL(2)$. 

Suppose that $\gamma ,\gamma ^{\prime }\in G({\Bbb{Q}})$. Following the
well-known definition, if there is a $g\in G(\overline{{\Bbb{Q}}})$ such that $%
\gamma ^{\prime }=g^{-1}\gamma g$ then we call $\gamma ,\gamma ^{\prime }$ 
\textbf{stably conjugate}. Here $\overline{{\Bbb{Q}}}$ denotes an algebraic
closure of ${\Bbb{Q}}$. As is well-known, two elements $\gamma ,\gamma
^{\prime }\in G({\Bbb{Q}})$ are conjugate under
$G(\overline{{\Bbb{Q}}})$ if and
only if they are conjugate under $G_1({\Bbb{Q}})$. See \cite{L3},
\cite{LL}.

Here is reformulation of this definition: If there is a
$g\in G_1({\Bbb{Q}})\hookrightarrow G_1({\Bbb{A}}_S)$ such
that $\gamma ^{\prime}=g^{-1}\gamma g$ then
$\gamma ,\gamma ^{\prime }$ are \textbf{stably conjugate}.
This condition has the following local analog which we make our
local definition:
Suppose that $\gamma ,\gamma ^{\prime }\in \overline{G(F)}$. If there
is a $g\in G_1(F\Bbb{)}$ such that 
$\gamma ^{\prime }=g^{-1}\gamma g$ then $\gamma ,\gamma ^{\prime }$ are 
\textbf{stably conjugate}.

\begin{remark}
We shall mention briefly that the paper of Springer-Steinberg
contains a comment (actually an exercise)
which may be worth mentioning in connection with
these definitions: \cite{SpSt}, ch IV, (2.15) says that if $E$ is
algebraically closed and $G\subset GL(n)$ is either a symplectic
group or an orthogonal group then then $g_1,g_2\in G(E)$ are
conjugate in $G(E)$ if and only if they are conjugate in $GL(n,E)$.
(I believe it is assumed throughout that $E$ is not of characteristic $2$.)
In particular, if $F$ is a p-adic field, $\overline{F}$ is its
algebraic closure and if $g_1,g_2\in G(F)$ are conjugate in
$GL(n,F)$ then they are conjugate in $G(\overline{F})$.
Recall the fact that if $g_1,g_2\in GL(n,F)$
are conjugate in $GL(n,\overline{F})$
then they are conjugate in $GL(n,{F})$ (``each conjugacy class
in $GL(n)$ is stable'', i.e., ``$GL(n)$ has no L-indistinguishability'').
These in conjunction imply that stable conjugacy in the
symplectic or orthogonal groups may be ``realized'' as
$GL(n,F)$-conjugacy.
\end{remark}

The global invariant trace formula is an equality of the form (9.53).
We will say that a distribution $T$ on
${\cal{H}}(\overline{G({\Bbb{A}}_S)})$ is 
\textbf{stably invariant} in case $T(f)=T(f^y)$ for all
$f\in {\cal{H}}(\overline{G({\Bbb{A}}_S)})$ and all
$y\in G_1({\Bbb{A}}_S)$. Not all the expressions
in the invariant trace formula (9.53) are stably invariant.
The object
of this section is to write the invariant trace formula
as the sum of stable
terms which we write in the form (as in \cite{L3})

\begin{equation}
\begin{array}{c}
{\rm tr}R_{cusp}(f)=
\sum_{\chi \in {\cal{X}}(G)}I'_\chi (f)\\
=-\sum_{\chi \in {\cal{X}}-{\cal{X}}(G)}I_\chi (f)
+S(f)
+\sum_{\mathcal{O}\in \mathcal{C}}SI_{\mathcal{O}}(f)
+{\rm\ ``endoscopic\ terms''}
\end{array}
\end{equation}
where $S=S_1+S_2+S_3$ is defined in (9.99-101),
$I'_\chi$ is equal to $I_\chi$ plus a term described in
\S 9.5.8 (g) below ($\chi\in
{\cal{X}}-{\cal{X}}(G)$), and
the endoscopic terms are a sum of terms of the form

\begin{equation}
\begin{array}{c}
{\rm\ ``endoscopic\ terms''}
=
{1\over 4}
\sum\Sb H\not= G\\ {\rm cuspidal}\endSb
\left[
\sum_{{\mathcal{O}}\in {\mathcal{C}}^H}SI^H_{\mathcal{O}}(f^H)
-\sum_{\theta\in E_H^G}<\theta,f^H>
\right]\\
={1\over 4}
\sum\Sb H\not= G\\ {\rm cuspidal}\endSb
\sum\Sb
\theta\in \Pi_u(H({\Bbb{Q}})\backslash H({\Bbb{A}}_S))\\
\theta\not= 1
\endSb
<\theta,f^H>
\end{array}
\end{equation}
which are stable distributions on endoscopic groups $H$ of $G$
(not $\overline{G({\Bbb{A}}_S)}$)
arising from the stable trace formula of $H$.
Here ${\cal{C}}^H=\{\{\gamma\}\ |\ \gamma\in H({\Bbb{Q}})\}$,
$E_H^G$ is the subset of
$\Pi_u(H({\Bbb{Q}})\backslash {H({\Bbb{A}}_S)})$,
and $f^H$ is a function on
$H({\Bbb{A}}_S)$ associated to
$f\in {\cal{H}}(\overline{G({\Bbb{A}}_S)})_{gen}$.
In the version presented below we will not
stabilize the $I_\chi$ in this paper.
We shall explain the terms ``$H$ a cuspidal endoscopic group of
$G$'' in (9.79) in the next subsection.

\subsection{Some basics on endoscopic groups}

If $T$ is an elliptic Cartan of $G=SL(2)$ defined over
a local or global field $F$ then let
\begin{equation}
D(T/F)=T_1(F)\backslash G_1(F)/G(F),
\end{equation}
where $G_1=GL(2)$ and $T_1=Cent(T,G_1)$.
We identify without further mention this double coset space with the
space
$\overline{T_1(F)}\backslash \overline{G_1(F)}/ \overline{G(F)}$,
where $\overline{G_1(F)}$ denotes the n-fold cover defined using the
cocyle $\beta$ in \S 2.1.
If $T$ is an elliptic Cartan of $SL(2)$ defined over
${\Bbb Q}$ then let
\begin{equation}
D(T/{\Bbb{A}})
=\oplus_p T_1({\Bbb{Q}}_p)\backslash G_1({\Bbb{Q}}_p)/G({\Bbb{Q}}_p).
\end{equation}
We embed $D(T/{\Bbb{Q}})$ into $D(T/{\Bbb{A}})$
or $D(T/{\Bbb{A}}_S)$ diagonally.

For general definitions, see \cite{L3}.
The expressions $SI_{{\cal{O}}_H}$ and $f^H$ are discussed there but
will also be mentioned (indirectly) below.

\subsection{Stabilizing the elliptic regular terms}

We shall call an element
$\gamma\in G({\Bbb{Q}})$ (resp.,
$(\gamma,\zeta) \in \overline{G(F)}$) \textbf{semisimple} if $\gamma $ is
diagonalizable and {\bf regular} if its eigenvalues
$\gamma_1,\gamma_2$ are distinct.
We shall call a semisimple element
$\gamma \in G({\Bbb{Q}})$
(resp., $(\gamma,\zeta) \in \overline{G(F)}$) \textbf{elliptic}
(or $G$-\textbf{elliptic}) if $\gamma $ belongs to
an elliptic Cartan.

We first shall recall the idea of stabilizing the elliptic regular terms of
(9.78) following \cite{LL}, \cite{L3}, \cite{Ko1}.

For $\gamma \in G({\Bbb{Q}})$, let
${\mathcal{O}}_\gamma \in {\mathcal{C}}$ denote
the equivalence class containing $\gamma $. 

We now turn to the trace formula in the rank $0$ case, as discussed in
chapter 3 above.

\begin{proposition} (Poisson summation)
Let $T\subset G$ be an elliptic Cartan defined over ${\Bbb{Q}}$. For any
$\phi \in C_c^\infty (\overline{T({\Bbb{A}}_S)})_{gen}$, we have 
\[
{\mathrm{vol}}(T({\Bbb{Q}})\backslash \overline{T({\Bbb{A}}_S)})
\sum_{\gamma \in T({\Bbb{Q}})}\phi (\gamma ) \\
=\sum_{\theta \in \Pi _u(T({\Bbb{Q}})\backslash \overline{T({\Bbb{A}}_S)})}
<\theta,\phi>.
\]
If $\theta =\otimes _p\theta _p$ and $\phi=\otimes _p \phi_p$ then 
\[
<\theta ,\phi>=\prod_p{\rm tr}\theta _p(\phi_p).
\]
\end{proposition}

\begin{remark}
({\bf A preview of coming attractions})
We shall eventually apply this with $\phi (\gamma )$ equal to

\begin{equation}
I_\gamma ^\kappa (f)=
[D(T/{\Bbb{Q}}): D(T/{\Bbb{A}}_S)]^{-1}
\sum_{\delta\in
D(T/{\Bbb{Q}})\backslash D(T/{\Bbb{A}}_S)}
\kappa (\det \delta)I_{\cal{O}}(f^\delta),
\end{equation}
where $\kappa \in \Pi _u(D(T/{\Bbb{Q}})\backslash D(T/{\Bbb{A}}_S))$,
$\gamma\in T({\Bbb{Q}})$ elliptic, and
$f\in {\cal H}(\overline{G({\Bbb{A}}_S)})$.
These terms (9.82) do not occur in the invariant trace formula.
However, after grouping terms as in (9.83) below and substituting
the expression in Lemma 328 below, all the terms (9.82),
$\gamma\in T({\Bbb{Q}})$, will occur in the trace formula
except for $\gamma =\pm 1$. The terms corresponding to
$\gamma =\pm 1$ will be discussed in section 9.5.5 below.
The terms (9.82) with $\kappa = 1$, $\gamma\not= \pm 1$,
will comprise part of the sum over
${\cal{O}}\in {\cal{C}}$ in (9.78). The
terms (9.82) with $\kappa \not= \pm 1$,
will comprise part of the sum in (9.79).  
\end{remark}

Next we record a corollary of the above proposition 
which we will substitute into the trace formula later.

\begin{corollary}
For each elliptic torus T, we have 
\[
vol[T({\Bbb{Q}})\setminus \overline{T({\Bbb{A}}_S)}]\cdot
\sum_{\gamma \in T({\Bbb Q})}
I_\gamma ^\kappa (f)=\sum_\theta <\theta ,I_{*}^\kappa (f)>
\]
where $I_{*}^\kappa $ is the unstable orbital integral defined in
(9.82)
above (for its properties, as a function on $T$, in particular on the
singular set of $T$, see \cite{LL} or section 9.5.5 below).
\end{corollary}

This expression will form part of what is meant by (9.79).

As we will see in Lemma 340 below,
we may use \cite{LL} to to relate $<\theta,I_{*}^\kappa (f)>$ to
a linear combination of characters of automorphic representations of
$\Pi_u(\overline{G({\Bbb{A}}_S)})$. If $\mathcal{O}=\mathcal{O}_\gamma $
and $\gamma \in \overline{T({\Bbb{A}}_S)}$ is elliptic regular, let 
\begin{equation}
\begin{array}{c}
I^{st}_{\cal{O}}(f)=\sum_{\delta\in D(T/{\Bbb{Q}})}
I_{\cal{O}}(f^\delta),\\
SI_{\cal{O}}(f)=I^1_{\gamma}(f)\; .
\end{array}
\end{equation}
These distributions $SI_{\cal{O}}$ are stable.

\begin{lemma}
(a)\ If $\mathcal{O}=\mathcal{O}_\gamma $ and
$\gamma \in \overline{T({\Bbb{A}}_S)}$ is
elliptic regular, we have 
\[
I^{st}_{\mathcal{O}}(f)
=\sum_{\kappa \in \Pi _u(D(T/{\Bbb Q})\backslash D(T/{\Bbb{A}}))}
I_\gamma ^\kappa (f)
\]
(b) If $\mathcal{O}=\mathcal{O}_\gamma $ and
$\gamma \in \overline{T({\Bbb{A}}_S)}$ is elliptic regular, we have 
\[
I_{h^{-1}\gamma h}^\kappa (f)=\kappa (\det h)I_\gamma ^\kappa (f)
=I_\gamma^\kappa (f^h),\ \ \ \ \ \ h\in \overline{G_1({\Bbb{A}}_S)}.
\]
\end{lemma}

Any function $g$ satisfying $g(h^{-1}\gamma h)=\kappa (\det (h))g(\gamma )$,
for all $h\in \overline{G_1({\Bbb{A}}_S)}$, will be called
\textbf{$\kappa $-stable}.

proof: It is elementary that if $A$ and $B$ are
discrete countable abelian groups, with $A\subset B$ of finite
index, and if $f$ is any compactly supported function on $B$ then 
\begin{equation}
\sum_{a\in A}f(a)=[B:A]^{-1}\sum_{\kappa \in (B/A)^{\wedge }}\sum_{b\in
B}\kappa (b)f(b),
\end{equation}
whenever the sums converge absolutely. We apply this to our
sum over $\delta $: 
\[
\sum_{\delta \in D(T/{\Bbb{Q}})}
I_{\cal{O}}(f^\delta)
=
\sum_{\kappa \in (D(T/{\Bbb{A}})/D(T/{\Bbb{Q}}))^{\wedge}}
[D(T/{\Bbb{A}}):D(T/{\Bbb{Q}})]^{-1}
\sum_{\delta \in D(T/{\Bbb{A}})}
\kappa (\delta )I_{\cal{O}}(f^\delta),
\]
which we may equate with
\begin{equation}
\sum_{\kappa \in \Pi_u(D(T/{\Bbb{A}})/D(T/{\Bbb{Q}}))}
I^\kappa_\gamma (f).
\end{equation}
$\Box$

\subsection{Stabilizing the $A$-elliptic regular terms}

We shall call a semisimple element
$\gamma \in G({\Bbb{Q}})$
(resp., $\gamma \in \overline{G(F)}$) $A$-\textbf{elliptic} regular if
$\gamma $ is conjugate to an element $a\in A(F)$
which is elliptic regular in $A(F)$ (i.e., $a\not= \pm 1$).

If $\mathcal{O}=\mathcal{O}_\gamma $ and
$\gamma \in G({\Bbb{Q}})$ is $A$-elliptic regular, let 
\begin{equation}
\begin{array}{c}
I^{st}_{\cal{O}}(f)=
\sum_{\delta\in D(A/{\Bbb{Q}})}
I_{\cal{O}} (f^\delta),\\
I^\kappa_\gamma(f)
=[D(A/{\Bbb{Q}}):D(A/{\Bbb{A}}_S)]^{-1}
\sum_{\delta\in D(A/{\Bbb{A}}_S)}
\kappa(\delta) I_{\cal{O}} (f^\delta),
\end{array}
\end{equation}
where
\[
\begin{array}{c}
D(A/{\Bbb{Q}})
=A({\Bbb{Q}})Z(G_1({\Bbb{Q}}))\backslash G_1({\Bbb{Q}})/G({\Bbb{Q}}) \\
D(A/{\Bbb{A}}_S)
=A({\Bbb{A}}_S)Z(\overline{G_1({\Bbb{A}}_S)})
\backslash \overline{G_1({\Bbb{A}}_S)}/G({\Bbb{A}}_S) \\
\kappa\in \Pi_u(D(A/{\Bbb{Q}})\backslash D(A/{\Bbb{A}}_S)).
\end{array}
\]
Let
\[
SI_{\cal{O}}(f)=I^1_{\gamma} (f),
\]
From (9.89), we have
\begin{equation}
I^{st}_{\cal{O}}(f)=
\sum_{\kappa\in \Pi_u(D(A/{\Bbb{Q}})\backslash D(A/{\Bbb{A}}_S))}
I^\kappa_\gamma(f)
\end{equation}
\begin{lemma}
If $\mathcal{O}=\mathcal{O}_\gamma $ and
$\gamma \in G({\Bbb{Q}})$ is $A$-elliptic regular, then
$SI_{\mathcal{O}}(f)$ is stable.
\end{lemma}

Recall that this term $SI_{\cal{O}}$ is defined using
\begin{equation}
I_{\cal{O}}(f)
=\sum_{\gamma\in {\cal{O}}}a^G(S,\gamma)I_G(\gamma,f)
+{1\over 2}\sum_{\gamma\in {\cal{O}}\cap A({\Bbb{Q}})}
a^A(S,\gamma)I_A(\gamma,f).
\end{equation}
The $I_M(\gamma,f)$ were defined by
\[
I_G(\gamma,f)=J_G(\gamma,f),
\]
which is an ``ordinary'' orbital integral over
$\overline{G({\Bbb{A}}_S)}_\gamma\backslash \overline{G({\Bbb{A}}_S)}$, and
\[
I_A(\gamma,f)=J_A(\gamma,f)-\hat{I}_A^A(\gamma,\phi_A(f)).
\]
(It turns out that
$\overline{G({\Bbb{A}}_S)}_\gamma\subset \overline{A({\Bbb{A}}_S)}$,
with equality if $\gamma$ is an $N$-th power, and
$J_G(\gamma,f)=0$ if $\gamma$ is not an $N$-th power - provided
$f$ is genuine as we have assumed.)
To see that the sum
$\sum_{\gamma\in {\cal{O}}}a^G(S,\gamma)I_G(\gamma,f)$
is stably invariant, we claim that $I_G(\gamma,f)$ is stably invariant.
This was proven in chapter 5.

The term $J_A(\gamma,f)$ is not invariant. We examine it
more closely in Lemma 332 and (9.91) below.
(See also \cite{A16}, Corollary 4.2 for a general formula for it
in the case $n=1$.)

The term $\hat{I}_A^A(\gamma,\phi_A(f))$ is not invariant. It is the
Fourier transform of the map
\[
\chi\rightarrow {\rm tr}(I_A(\chi,f)R_A(\chi,B)),
\]
see (9.67). The instability of $f\longmapsto \phi_A(f)$ is
examined in Lemma 341 above.

Let $F={\Bbb{Q}}_p$, $p<\infty$, $K=G({\Bbb{Z}}_p)$ which we
identify with a {\em subset} $(K,1)$ of $\overline{G(F)}$,
$\sigma\in \Pi(\overline{A(F)})$. Let $K_{I_A(\sigma)}(k,k')$
denote the kernel function of the induced representation
$I_A(\sigma)$ on the space of functions on $K$, so
\begin{equation}
{\rm tr}(I_A(\sigma,f))
=\int_K K_{I_A(\sigma)}(k,k)dk .
\end{equation}

\begin{lemma}
If $F(\sigma)={\rm tr}(I_A(\sigma,f))$ then
\[
F^\vee(\gamma)
=|a_1-a_2|J_G(\gamma,f),
\]
for $\gamma =(a,\zeta)\in \overline{A(F)}$, $a\not= \pm 1$.
\end{lemma}

proof: This was proven in chapter 5. $\Box$

As a corollary, it follows from this lemma without to much work
that $f\longmapsto I_G(\gamma,f)$ is stable, for
$\gamma =(a,\zeta)\in \overline{A(F)},\ a\not= \pm 1$.

Let $\delta(g)=|a_1/a_2|$, if $g=(nak,\zeta)\in \overline{G_1(F)}$. Let
$N(g)$ denote the operator on the space of functions on $K$ given by
\[
f(k)\longmapsto f(k)\log\delta(kg).
\]
This operator also occurs in connection with the terms in the
invariant trace formula associated to the normalized intertwining
operators.

\begin{lemma}
For $g\in G_1(F)$, we have
\[
{\rm tr}(I_A(\sigma,f)N(g))
=\int_K K_{I_A(\sigma)}(k,k)\log\delta(kg)dk.
\]
\end{lemma}

\begin{lemma}
For $g\in G_1(F)$, if
$F(\sigma)={\rm tr}(I_A(\sigma,f)N(g))$ then
\[
J_A(a,f^{g^{-1}})-J_A(a,f)
=F^\vee(a)+F^\vee(a^w)
\]
\end{lemma}
proof: This was proven in chapter 5. $\Box$

If $v>2$ then we may represent $D(A/F)$ by
\[
\{
\left(
\begin{array}{cc}
a&0\\
0&1
\end{array}
\right)\ |\ a=1,\epsilon, \pi, \epsilon\pi \},
\]
where $\pi$ is a uniformizer and $\epsilon$ is a unit
which is not a quadratic residue. If $v=2$ the situation is similar...
Therefore, if $\kappa$ is non-trivial
we may apply Lemma 332 above to compute
\[
\sum_{\delta\in D(A/F)}
\kappa(\delta)J_A(\gamma,f^\delta).
\]

If $x=a
\left(
\begin{array}{cc}
1&n\\
0&1
\end{array}
\right) k$, $n\in {\Bbb{A}}_S$, then
\begin{equation}
v_A(x)
=v(x,0)=\sum_v v(0,x_v)=
\sum_v \log\sqrt{|1+n_v^2|_v}
\end{equation}
and
\begin{equation}
\begin{array}{c}
J_A(\gamma,f)
=|D(\gamma)|^{1/2}
\int_{\overline{G({\Bbb{A}}_S)}_\gamma
\backslash \overline{G({\Bbb{A}}_S)}}
f(x^{-1}\gamma x)v_A(x)dx\\
=
\sum_w\prod_v
|D(\gamma_v)|_v^{1/2}
\int_{\overline{G({\Bbb{Q}}_p)}_\gamma
\backslash \overline{G({\Bbb{Q}}_p)}}
f_v(x_v^{-1}\gamma_v x_v)v_A(x_w)dx_v\\
=\sum_v |D(\gamma_v)|_v^{1/2} J_A(\gamma_v,f_v)
\prod_{w\not= v}|D(\gamma_w)|_w^{1/2} J_G(\gamma_w,f_w).
\end{array}
\end{equation}

\begin{remark} This comprises metaplectic analog of
the $a\not= \pm 1$
part of the distribution denoted $S_2$ in \cite{LL}. It
will be combined with other distribution
described in the next section to yield a stably
invariant distribution.
\end{remark}

\subsection{Stabilizing the singular terms}

A term $I_{\cal{O}}(f)$ is called {\bf singular} if ${\cal{O}}$
does not contain any $M$-elliptic elements, for any $M\in \{A,G\}$.
An element $\gamma\in G({\Bbb{Q}})$ is called {\bf singular}
if its eigenvalues $\gamma_1,\gamma_2$ are equal (i.e.,
$\gamma$ is not regular).

In this subsection, our goal is to show that part of
the singular terms of the geometric side of the invariant
trace formula (9.53) ``combine'' with the sum of the terms
$I_\gamma ^\kappa (f)$,
with $\kappa \neq 1$ and with
$\gamma \in T({\Bbb{Q}})$ elliptic \emph{singular}
(i.e., $\gamma =\pm 1$) to contribute to (9.79). 
If ${\cal O}=\{z\}$, $z=\pm 1$, then define
\begin{equation}
SI_{\cal O}(f)=c\cdot f(z),
\end{equation}
where $c=vol(G({\Bbb Q})\backslash \overline{G({\Bbb{A}}_S)})$ is a volume.

\begin{lemma} If $z\in {\cal{O}}$, with $z=\pm 1$, then
\[
I_{\cal O}(f)-SI_{\cal O}(f)=
\sum\Sb \gamma\in {\cal{O}}\\\gamma\not= z\endSb
a^G(S,\gamma)I_G(\gamma,f)+
{1\over 2}a^A(S,z)I_A(z,f).
\]
\end{lemma}
proof: This equation follows from (9.88). $\Box$

\begin{definition}
Let $\kappa$ be a character of
$D(T/F)$ of order $2$.
Define $\Phi^{T/\kappa}$ on the regular set of $\gamma\in T(F)$ by
\[
\begin{array}{c}
\Phi^{T/\kappa}(\gamma,f)=
\Delta(\gamma,T,\kappa)
\int_{\overline{G_1(F)}_\gamma\backslash \overline{G_1(F)}}
\kappa(\det{g})f(g^{-1}\gamma g)dg\\
=
\lambda(E/F,\psi)
\kappa({\gamma_{1,v}-\gamma_{2,v}\over \gamma_{1,v}^0-\gamma_{2,v}^0})
\int_{\cal{O}}\epsilon(x)f(x)|\omega_\gamma|,
\end{array}
\]
where $\epsilon(g^{-1}\gamma g)=\kappa(\det(g))$,
$\gamma^0\in \overline{T(F)}$
is fixed, and $\cal{O}$ denotes the conjugacy class of
$\gamma$ in $\overline{G(F)}$.
\end{definition}

We leave the relation between the global expressions
$\Phi^{T/\kappa}(\gamma,f)$ and $I^\kappa_\gamma(f)$ to
the reader (use the Hilbert product formula).

Let ${\cal{O}}(\gamma)$ denote the orbit of
$\gamma$ in $\overline{G(F)}$ and,
for $a\in F^\times$, which we identify with the center of $G_1(F)$, let
\[
n(a)=a
\left(
\begin{array}{cc}
1&1\\ 0&1
\end{array}
\right).
\]

\begin{definition} If $z\in Z(\overline{G(F)})$, then we define
\[
\begin{array}{c}
\Phi^{T/\kappa}(z,f)\\
=
\left\{
\begin{array}{c}
\int_{{\cal{O}}(n(z))}
f(x)dx,\ \ \ \ T=A,\ \kappa =1\\
\lambda(E/F,\psi)
\kappa({\pi_E-\overline{\pi_E}\over \gamma_{1,v}^0-\gamma_{2,v}^0})
\int_{{\cal{O}}(n(z))}\epsilon(x)f(x)dx,\ \ \ \ T_1(F)=E^\times,
\kappa\in \Pi_u(F^\times),
\end{array}
\right.
\end{array}
\]
where $\pi_E$ is a uniformizer for $E/F$, $T_1(F)=E^{\times 1}$,
and
\[
\epsilon(g^{-1}n(z)g)=
\left\{
\begin{array}{c}
\kappa(-b),\ \ \ \ b\not= 0,\\
\kappa(c),\ \ \ \ b= 0,\\
\end{array}
\right.
\]
if
\[
g^{-1}n(z)g=
\left(
\begin{array}{cc}
a&b\\c&d
\end{array}
\right).
\]
\end{definition}

\begin{lemma} If $\kappa\in \Pi_u(F^\times)$ is the character
of order 2
associated to $E/F$ be local class field theory ($\kappa =1$
if $E=F\times F$) and $T_1(F)=E^\times$, then
$\gamma\longmapsto \Phi^{T/\kappa}(\gamma,f)$ is a smooth function on
$T(F)$.
\end{lemma}

proof: This follows from the results on Shalika germs. It is proven
in chapter 5. $\Box$

\begin{lemma} For $z\in Z(\overline{G(F)})$, we have
\[
\begin{array}{c}
\lim_{\gamma\rightarrow z}
\int_{\overline{G_1(F)}_\gamma\backslash \overline{G_1(F)}}
\kappa(\det{g})f(g^{-1}\gamma g)dg\\
=c{1\over L(1,\kappa)}
\int_{\overline{G_1(F)}_{n(1)}\backslash \overline{G_1(F)}}
\kappa(\det{g})f(g^{-1}n(z) g)dg
\end{array}
\]
where $c=|\gamma_{1,v}^0-\gamma_{2,v}^0|$.
\end{lemma} 

proof: This was proven in chapter 5. $\Box$

\begin{lemma} Let $\kappa$ be a character of
${\Bbb{A}}_S^\times $ of order 2 which we extend to
$\overline{G_1({\Bbb{A}}_S)}$ via the $\det$ map,
and $T$ be a Cartan of $G$ associated to $\kappa$ by class field
theory then
\begin{equation}
\sum_{\delta\in D(A/{\Bbb{A}}_S)}
\kappa(\delta)I_G(n(z),f^\delta)
=c
\lim\Sb \gamma\rightarrow z \\ \gamma\in T({\Bbb{Q}})\endSb
\sum_{\delta\in D(T/{\Bbb{A}}_S)}
\kappa(\delta)I_G (\gamma, f^\delta),
\end{equation}
for $z =\pm 1$, $\kappa \in \Pi _u(G_1({\Bbb{A}}_S))$ of order 2,
where the constant $c=c(T,\kappa,S)$ is given by
\[
c=\prod_v |\gamma_{1,v}^0-\gamma_{2,v}^0|_v^{-1}L(1,\kappa_v) .
\]
\end{lemma}

\begin{remark} The terms on the right hand side of the above
lemma go into (9.97) below. To see this, one relates them to
$I_{\cal{O}}(f),\ z\in {\cal{O}}$, by using (9.94) below.

The terms on the left hand side will examined following the proof
below.
\end{remark}

proof: This is is a consequence of the previous lemma and the
definitions. $\Box$

If $z\in {\cal{O}}$ is an equivalence class of $G({\Bbb{Q}})$
as above then part of $I_{\cal{O}}(f)$ may be written as
\begin{equation}
\begin{array}{c}
\sum_{\gamma\in {\cal{O}}} J_G(\gamma,f)
=J_G(z,f)+\sum_{a\in {\Bbb{Q}}^{\times}/{\Bbb{Q}}^{\times 2}}
\int_{\overline{G({\Bbb{A}}_S)}_{n(1)}
\backslash \overline{G({\Bbb{A}}_S)}} f(g^{-1} z\cdot
\left(
\begin{array}{cc}
1 & a\\
0 & 1
\end{array}
\right)
g)dg\\
=J_G(z,f)+
\sum_{\kappa \in \Pi_u({\Bbb{A}}_S^{\times}/{\Bbb{A}}_S^{\times 2})}
\sum_{a \in {\Bbb{A}}_S^{\times}/{\Bbb{A}}_S^{\times 2}}
\kappa (a)
\int_{\overline{G({\Bbb{A}}_S)}_{n(1)}
\backslash \overline{G({\Bbb{A}}_S)}} f(g^{-1} z\cdot
\left(
\begin{array}{cc}
1 & a\\
0 & 1
\end{array}
\right)
g)dg\\
=J_G(z,f)+
\sum_{\kappa \in \Pi_u({\Bbb{A}}_S^{\times}/{\Bbb{A}}_S^{\times 2})}
\int_{\overline{G_1({\Bbb{A}}_S)}_{n(1)}
\backslash \overline{G_1({\Bbb{A}}_S)}} \kappa(\det{g})
f(g^{-1} n(z) g)dg,
\end{array}
\end{equation}
which we write as $J_G(z,f)$ plus the sum over
$\kappa \in \Pi_u({\Bbb{A}}_S^{\times}/{\Bbb{A}}_S^{\times 2})$
of the left hand side of the equation in Lemma 338 above.

Regarding the expression
\[
{a^A(S,z)\over 2}
\sum_\delta \kappa(\delta)
(J_A(z,f^\delta)-\hat{I}_A^A(z,\phi_A(f^\delta)))),
\]
we must also compute $a^A(S,z)$ and discuss the terms
$\sum_\delta \kappa(\delta)J_A(z,\phi_A(f^\delta))$ and
$\sum_\delta \kappa(\delta)\hat{I}_A^A(z,\phi_A(f^\delta))$.

Computation of $a^A(S,z)$: By Proposition 279,
\[
a^A(S,z)=vol(A({\Bbb{Q}})\backslash \overline{A({\Bbb{A}}_S)^1}).
\]

Computation of
$\sum_\delta \kappa(\delta)J_A(z,\phi_A(f^\delta))$:
See Example 275 above. These terms
comprise the $a=\pm 1$ part of the distribution denoted by
$S_2(f)$ ``missing'' from (9.94).

Computation of $\sum_\delta \kappa(\delta)\hat{I}_A^A(z,\phi_A(f^\delta))$:
First, we note that the map
$\phi_A:{\cal{H}}(\overline{G({\Bbb{A}}_S)}) \rightarrow I(A({\Bbb{A}}_S))$
of section 9.4.2 is a ``weighted
character'' of an induced representation:
\begin{equation}
\begin{array}{c}
\phi_A(\sigma,f)=
{\rm tr}(R_{\overline{B}|B}(\sigma_\lambda)^{-1}
R'_{\overline{B}|B}(\sigma_\lambda)I_A(\sigma_\lambda,f))|_{\lambda =0}\\
=\sum_v
{\rm tr}(R_{\overline{B}|B}(\sigma_v)^{-1}
R'_{\overline{B}|B}(\sigma_v)I_A(\sigma_v,f_v))
\prod_{w\not= v}{\rm tr}(I_A(\sigma_w,f_w)).
\end{array}
\end{equation}
for $\sigma\in \Pi_u(\overline{A({\Bbb{A}}_S)})$.
The inverse Fourier transform of the function
$\sigma\longmapsto \phi_A(\sigma, f)$ evaluated at $a=1$ is
essentially the distribution denoted by $S_1(f)$ in \cite{LL}.

\subsection{The geometric side of the stable(r) trace formula}

The stable elliptic regular terms are 
\begin{equation}
\sum\Sb {\mathcal{O}}\in {\mathcal{C}}_G \\
\mathcal{O}\;\mathrm{regular\ elliptic} \endSb
SI_{\mathcal{O}}(f),
\end{equation}
where the $SI$'s are defined by (9.83), 

The following definition, remark, and two lemmas will not be needed here.
There are used to introduce notation and to explain (9.98) below.
We do this more for the sake of comparing the difference
between the ``endoscopic'' terms occuring here and in \cite{LL}, \S 5.

\begin{definition} If $\pi$ is an irreducible admissible
genuine representation of
$\overline{G({\Bbb Q}_p)}$ then let
$L(\pi)$ denote the set of inequivalent representations
$\pi^g$ where $g\in \overline{G_1({\Bbb{Q}}_p)}$.
\end{definition}

\begin{remark} For $\pi$ as above and $n>1$ then
\[
\# L(\pi)=
\prod_{v|p}\#( k_v^\times/k_v^{\times 2}).
\]
\end{remark}

If $\theta_1$ is a character
of ${T_1({\Bbb Q}_p)}$ and $\pi_1(\theta_1)$ is the
irreducible admissible representation of
${G_1({\Bbb Q}_p)}$ associated to it, let
$\theta =\theta_1|_{{T({\Bbb{Q}}_p)}}$ and let
$\Pi(\theta)$ denote the set of the irreducible
constituents of
$\pi_1(\theta_1)|_{G({\Bbb{Q}}_p)}$.
The following lemma will not be used here. Note that it related
the distribution
\[
\begin{array}{c}
C_c^\infty(\overline{G({\Bbb{A}}_S)})_{gen}
\rightarrow
C_c^\infty(\overline{T({\Bbb{A}}_S)})_{gen}\\
f\longmapsto I_*^\kappa(f)
\end{array}
\]
to a map
\[
\begin{array}{c}
C_c^\infty(T({\Bbb{A}}_S))'
\rightarrow
C_c^\infty(\overline{G({\Bbb{A}}_S)})'_{gen}\\
\theta \longmapsto
\prod_{p\in S}\sum_{\pi _p\in \Pi (\theta_p)}
\kappa _{*}(\pi _p)\chi _{\pi _p}(f_p),
\end{array}
\]
after relating $C_c^\infty(\overline{T({\Bbb{A}}_S)})_{gen}$
to $C_c^\infty(T({\Bbb{A}}_S))$.
\begin{lemma}
For each elliptic torus T and each factorizable
$\overline{f}\in {\cal{H}}(\overline{G({\Bbb{A}}_S)})_{gen}$, we have 
\[
\begin{array}{c}
<\theta ,I_{*}^\kappa (\overline{f} )> \\ 
=[D(T/{\Bbb{Q}}):D(T/{\Bbb{A}}_S)]^{-1}
\prod_{p\in S}
\sum_{\delta\in D(T/{\Bbb Q}_p)}
\sum_{\pi _p\in \Pi (\theta_p^\delta)}
\kappa _{*}(\pi _p)\chi _{\pi _p}(f_p),
\end{array}
\]
where $\kappa _{*}(\pi _p):=\kappa (\delta)$ if
$\pi _p\in \Pi (\theta_p^\delta)$,
${\theta }_p^h$ denotes the conjugate of ${\theta }_p$ by $h$,
${\theta }_p$ is a unitary genuine irreducible
representation of $\overline{T({\Bbb{A}}_S)}$, and
$f_p(g)=\overline{f}_p(g,1)$.
\end{lemma}

proof: This is proven in chapter 5. $\Box$
\begin{remark}
This result is essentially the metaplectic analog of
\cite{LL}, (2.4-5), p. 762. This is, after a suitable 
choice of global measures, related to a special case of
\cite{J}, (3.42) which was only proven under in the case $n$ odd.
\end{remark}

By this lemma and the trace formula for rank 0 (in
Proposition 326)
the terms 
\begin{equation}
\sum\Sb T \; \\ G-\mathrm{elliptic} \endSb
{\rm vol}(T({\Bbb{Q}})\backslash \overline{T({\Bbb{A}}_S)})
\sum_{\gamma \in T({\Bbb{Q}})}
\sum\Sb \kappa \in \Pi _u(D(T/{\Bbb Q})\backslash D(T/{\Bbb{A}}_S))
\\ \kappa \neq 1
\endSb
I_\gamma ^\kappa (f)
\end{equation}
may be expressed as a linear combination of characters of representations 
\begin{equation}
\sum_T
\sum_\theta c_\theta tr(\pi (\theta )(f)),
\end{equation}
for some constants $c_\theta$, where $\theta$ runs over
$\Pi_u(T({\Bbb{Q}})\backslash \overline{T({\Bbb{A}}_S)})$.
Note that we may rewrite the terms in (9.97) as
\[
\sum_{H\ cuspidal}
\sum\Sb {\cal O}_H\not= \pm 1\\
a\ conj\\
class\ in\ H
\endSb
SI_{{\cal O}_H}^H(f^H)
\]
in the notation of (9.79).

In (9.97) the terms $\gamma =\pm 1$ occuring
in the middle sum come from the singular elliptic terms in
Corollary 327 above.

\subsection{Stabilizing the remaining terms}

Let $F={\Bbb{Q}}_p$.

\begin{lemma} For all $g\in \overline{G_1(F)}$,
\[
\begin{array}{c}
{\rm tr}(R_{\overline{B}|B}(\sigma)^{-1}
R'_{\overline{B}|B}(\sigma)I_A(\sigma,f^{g^{-1}}))
-{\rm tr}(R_{\overline{B}|B}(\sigma)^{-1}
R'_{\overline{B}|B}(\sigma)I_A(\sigma,f))\\
={\rm tr}(I_A(\sigma,f)N(g))+{\rm tr}(I_A(\sigma^w,f)N(g))
\end{array}
\]
\end{lemma}

proof: This is proven in chapter 5. $\Box$

Let
\begin{equation}
\begin{array}{c}
S_1(f)\\
=\sum_{\chi\in{\cal{X}}-{\cal{X}}(G)}
\sum_v {1\over 2}
\int_{\Pi(A,\chi)}
{\rm tr}(R_{\overline{B}|B}(\sigma_v)^{-1}
R'_{\overline{B}|B}(\sigma_v)I_A(\sigma_v,f_v))
\prod_{w\not= v}
{\rm tr}(I_A(\sigma_w,f_w))|d\pi |.
\end{array}
\end{equation}

Let
\begin{equation}
S_2(f)=
-{1\over 2}
\sum_v \sum_{\{a,a^w\}\in A({\Bbb{Q}})}
{|D(a)|^{1/2}_v
J_A(a,f_v)\over L(1,{\Bbb{Q}}_v)}
\prod_{w\not= v}
{|D(a)|^{1/2}_w J_G(a,f_w)\over L(1,{\Bbb{Q}}_w)}.
\end{equation}

Let
\begin{equation}
\begin{array}{c}
S_3(f)=\\
=Res_{s=0}{1\over s}
\sum_{z\in Z(G({\Bbb{Q}}))}
\sum\Sb
x\in {\Bbb{Q}}^\times /{\Bbb{Q}}^{\times 2}
x\not= 1 \endSb
vol(Z(G({\Bbb{Q}}))\backslash Z(\overline{G({\Bbb{A}}_S)}))\times \\
\times
\int_{\overline{G({\Bbb{A}}_S)}_{n(1)}
\backslash \overline{G({\Bbb{A}}_S)}}
f(g^{-1}z
\left(
\begin{array}{cc}
1&x\\ 0&1
\end{array}
\right)
 g)\lambda(g)^{-s}dg ,
\end{array}
\end{equation}
where $\lambda(g)=\delta(g)\delta(wg)$.

The following will not be needed. It is the analog of a result in
\cite{LL}, the proof of which used the analog of Lemma 346 above.
\begin{conjecture}
Let $S=S_1+S_2+S_3$. Then $S$ is stable: 
$S(f^g)=S(f)$ for all $g\in \overline{G_1({\Bbb{A}}_S)}$.
\end{conjecture}

Perhaps the proof of this (assuming it is correct)
would likely be similar to the proof of its analog in \cite{LL}.

\subsection{The stable(r) trace formula}

This is the equality 
\begin{equation}
\sum_{\mathcal{O}\in \mathcal{C}}SI_{\mathcal{O}}(f)
=\sum_{\chi \in {\mathcal{X}}}I_\chi (f)
-\sum_{H\ cuspidal}\sum_{{\cal O}_H}SI_{{\cal O}_H}^H(f^H),
\end{equation}
where the $SI_{\mathcal{O}}$ are as in sections 9.5.3-9.5.5.

The geometric side of the stable(r)
trace formula is the sum of the following terms:

(a) the sum over $(G,S)$-classes ${\cal{O}}\subset G({\Bbb{Q}})$
such that ${\cal{O}}\cap Z(G({\Bbb{Q}}))\not= \emptyset$ of
\begin{equation}
SI_{\cal{O}}(f)=
\sum_{z\in ({\cal{O}}\cap A({\Bbb{Q}}))_{(A,S)}}
a^G(S,z)f(z),
\end{equation}
where $a^G$ is (a volume) is given in section 9.2.8 
(these do not include the ``nilpotent terms'', which
got combined with some other terms ...),

(b) the sum over elliptic regular classes 
${\cal{O}}\subset G({\Bbb{Q}})$
of
\begin{equation}
SI_{\cal{O}}(f)=
\sum_{\gamma\in {\cal{O}}_{(G,S)}}
\sum_{\delta\in D(T/{\Bbb{A}}_S)}
a^G(S,\gamma)J_G(\gamma,f^\delta),
\end{equation}
where $\gamma\in {\cal{O}}$, $T=G_\gamma$, and
$J_G$ (an orbital integral) and $a^G$ (a volume)
are given in section 9.2.8,

(c)  the sum over $A$-elliptic regular classes
${\cal{O}}\subset G({\Bbb{Q}})$
of
\begin{equation}
SI_{\cal{O}}(f)=
\sum_{\delta\in D(A/{\Bbb{A}}_S)}
I_{\cal{O}}(f^\delta),
\end{equation}

(d) plus terms associated to the geometric side of the 
trace formula for each cuspidal endoscopic of $G$,

(e) the ``geometric part'' of the endoscopic terms (9.79).

The spectral side of the stable(r) metaplectic
trace formula is the sum of the following terms:

(f) the sum over classes $\chi\in {\cal{X}}(G)$ of
\begin{equation}
I_\chi(f)=
{\rm tr}(\pi(f)),\ \ \ (\pi,G)\in\chi.
\end{equation}

(g) The sum over classes $\chi\in {\cal{X}}-{\cal{X}}(G)$
of
\begin{equation}
I_\chi(f)=-\sum\Sb \sigma\\
(\sigma,A)\in\chi\endSb
{m'(\sigma)\over m(\sigma)}
{\rm tr}I_A(\sigma,f).
\end{equation}

(h) the sum of the distributions $S_1$, $S_2$, and $S_3$

(i) the ``spectral part'' of the endoscopic terms (9.79)
(which we have omitted ...).

\section{Some food for thought}

It would be nice if we had

(1) provided a natural interpretation of the
stabilization of the terms in (h),

(2) uncovered endoscopic analogs of Arthur's
inductive formula of the invariant
distributions $I_{\cal{O}}$ and $I_\chi$ (if one exists),

(3) a stabilization of the terms in (f).


